\section{Формальная верификация}

\tab Формальная верификация относится к статическим методам верификации, т.е. проверка корректности устройства производится без фактического запуска RTL-кода. Вместо этого по отношению к модели применяются формальные (математические) методы, с помощью которых доказывается соответствие/несоответствие заданным в спецификации требованиям. 

\subsection{Базовые принципы работы}

Основой формальной верификации является алгоритм проверки модели (model checking algorithm). Этот алгоритм принимает на вход модель, которую мы проверяем, и формальную спецификацию, а на выход подает либо утверждение, что модель соответствует спецификации, либо контрпример этого утверждения в случае провала.

Кратко опишем алгоритм проверки модели при условии, что мы имеем формальное свойство $\phi$ и RTL-модуль $M$: 
\begin{enumerate}
    \item Создаем проверяющий автомат для $\neg \phi$.
    \item Извлекаем из $M$ конечный автомат $J$.
    \item Считаем умножение J и $\phi$.
    \item Если умножение непусто и содержит какой-то путь, то мы установили, что M содержит путь, который удовлтворяет не фи, и, следовательно, оповергает фи. Таким образом, можно сообщить о наличии ошибки, а в качестве контрпримера привести найденный путь.
    \item Если уможение пусто, то тогда считаем, что M удовлтворяет фи.
\end{enumerate}

Как было сказано ранее, формальная верификация крайне чувствительна к размеру и сложности проверяемого блока. На первый вгляд может показаться, что это вызвано сложностью алгоритма, но основная проблема заключается в размере извлеченного из модели конечного автомата. Умножение вместе проверкой на пустоту линейно относительно размера FSM и проверяющего автомата. При это размер автомата особой сложности не вносит, т.к. свойства, разрабатываемые верификиатором, как правило имеют небольшой размер. 

КА с другой стороны, может расти очень быстро в зависимости от некоторых особенностей модели. Рассмотрим их по очереди: 
\begin{itemize}
    \item Последовательная логика. Дает сильное усложенение для unbounded доказательств. За счет связи нынешнего состояния с предыдущим размер BDD (Binary Decision Diagram) увеличивается, но, что важнее, растет количество шагов индукции. Например, для FIFO шириной W и глубиной D количество соcтояний после reset – 2\^(W*D), а для unbounded доказательства понадобится как минимум D шагов индукции. Лучше всего для формальной верификации подходят блоки с глубиной не более 100.
    \item Ширина данных. Как правило, транспортировка данных подразумевает наличие широких шин и/или наличия в блоке банка данных (памяти). Все это напрямую влияет на увеличение кол-ва состояний системы. 
    \item Трансформация данных. Вывод из предыдущего пункта + дополнительное увеличение сложности за счет вычислений.
    \item Параметры. Т.к. для каждого набора параметров нужен свой прогон утверждений, то суммарное количество прогонов растет (особенно если параметр – это не бинарная величина). Более того, значение параметра может влиять на время прохождения доказательств (если система вовсе сойдется), например, в случае FIFO с параметризированной глубиной (см. выше).
\end{itemize}

Иначе говоря, попытка формально верифицировать почти любую модель средней сложности ведет к комбинаторному взрыву.

\subsection{Проблема комбинаторного взрыва}

Полностью решить проблему комбинаторного взрыва невозможно: в худшем случае сложность извлеченного из модуля конечного автомата будет расти экспоненциально от сложности самого модуля. Тем не менее, существует ряд решений, которые существенно упрощают представление КА. Расмотрим наиболее популярные сейчас. 

\subsubsection{Бинарная диаграмма решений}

BDD (Binary Decision Diagram) --- это компактное каноническое представление булевых функций. Каноническим называют такое представление, что для любых двух функционально эквивалентных функций оно будет одинаковым.  Introduction to VLSI Design Flow

Именно использование BDD вывело формальную верификацию в ряды практически применимых методов. Первые исследования показывали, что формальную верификацию можно применять для систем, содержащих вплоть до 10\^20 состояний, что на тот момент считалось недостижимым с использованием более классических методов. 

Для BDD фундаментальной является теорма разложения Шэннона, которая гласит, что любую булеву функцию $f(x_1, x_2, ..., x_k)$ от переменных ... мы можем представить в виде $f() = $. Последовательно применяя эту теорему к логике, описавающей проверяемую модель, можем получить конечный автомат в виде дерева решений. Каждый узел этого дерева представляет булеву функцию поддерева этого узла. 

Последовательность, в которой мы выбираем переменные для разложения Шэннона, называется порядком переменных. Прим этом, такое дерево решений называется упорядоченной, то есть различные переменные появляются в одном и том же порядке во всех путях от корневого узла графа до одной из терминальных вершин. 

Некоторые узлы дерева решений могут выражать одну и ту же функцию. Пользуясь ацикличностью графа, можем изпользовать два правила для упрощения дерева:
1. Если два узла представляют одну и ту же функцию, то есть, они изоморфны, то производим их слияние
2. Если у узла два изоморфных потомка, то удаляем узел 
Таким образом мы получем сокращенное упорядоченное ..., или же СУБДР (ROBDD). Чаще всего под термином БДР поздразумевают именно эту форму. 

Основным ограничением в использовании BDD является размер. Хоть BDD и является более компактной формой записи булевых функций, она не устраняет проблему комбинаторного взрыва, а лишь смягчает ее. К тому же, размер может сильно зависеть от порядка выбора переменных, в худшем случае эксопненциально. Причем алгоритмы поиска наилучшего упорядочивания сами по себе могут быть сложными. 

\subsubsection{Задача выполнимости булевых формул}

Как было показано ранее, доказательство на основе BDD может затрачивать огромное количество ресурсов. В то же время, для нахождения контпримера к свойству нам необязательно выражать КА модели полностью. Обычно мы ожидаем свойство провалиться в течение n первых тактов, поэтому нам достаточно просто развернуть КА до этого предела.

В таком случае проблему нахждения контрпримера можно свески к SAT-задаче. Действительно, пусть у нас есть первоначальное состояние I и свойство P. Введем свойство фи = не P, а также функцию смены состояний C. Наша задача --- установить, можем ли мы найти след длины n такой, что удовлетворяется фи. 

Пусть n = 2. Создадим свойство-свидетеля фи, то есть свойство, описывающее все возможные пути длины 2, при которых удовлетворяется фи. Тогда нам нужно решить сат проблемы для следующей формулы:


Далее пусть эн равно 3. Снова, мы используем ...
    
\subsection{Практическая проверка применимости}

Одной из целей данной работы является проверка применимости различных методов верификации. В отличие от симуляции, формальная верификация крайне чувствительна к размерам проверяемой системы, поэтому поставленная задача почти полностью сводится к практическому изучению максимальной сложности модели, с которой может работать инструмент формальной верификации.

Разумеется, однозначно определить предельный размер модуля невозможно. Тривиальные свойства, наподобие проверки valid-ready handshake, будут быстро доказываться и для очень сложных систем за счет маленького конуса влияния. Также на скорость прохождения доказательства влияют оптимизации, используемые в инструменте формальной верификации, сложность свойства, вычислительная мощность и др. Тем не менее, мы можем попробовать рассмотреть какую-то типовую ситуацию.

Так как найти контрпример или доказать свойство ограниченно значительно легче, чем доказать его полностью с построением BDD, то поставим две задачи: 1) выяснить максимальный предел состояний, с которым инструмент формальной верификации может выдать однозначный, неограниченный ответ и 2) выяснить зависимость числа тактов, для который справедливо свойство, от количеств состояний системы. 

В качестве модуля для проверки была выбрана 1r1w память. В этом модуле практически отсутствует дополнительная логика, и основной нагрузкой является массив данных. В качестве пробного свойства было выбрано свойство read-after-write consistency, записанное с использованием сложных для проверки операторов throughout и [->1]. 

\begin{lstlisting}
property p_rdy;
    logic [DATA_WIDTH-1:0] v_data;
    logic [ADDR_WIDTH-1:0] v_addr;
    @(posedge clk)
    ((i_wr_en, v_data = i_wr_data, v_addr = i_wr_addr) ##1 
        !wr_en throughout (i_rd_addr == v_addr)[->1])
    |-> o_rd_data == v_data;
endproperty
\end{lstlisting}

Для первой подзадачи получили следующие результаты: 

\begin{table}[H]
\centering
\begin{tabular}{|c|c||c|}
\hline
Сложность блока, flops & Сложность блока, gates & Время, с \\ \hline
147                    & 541                    & 0,1      \\ \hline
283                    & 993                    & 0,2      \\ \hline
555                    & 1897                   & 0,9      \\ \hline
1099                   & 3705                   & 2,1      \\ \hline
2187                   & 7321                   & 4,5      \\ \hline
4363                   & 14553                  & 35,2     \\ \hline
8715                   & 29017                  & 165,9    \\ \hline
17419                  & 57945                  & 729,4    \\ \hline
26123                  & 86873                  & 1206,6   \\ \hline
34827                  & 115801                 & 3905,9   \\ \hline
43531                  & 144729                 & 8360,2   \\ \hline
52235	               & 173657			        & 10158,2  \\ \hline


8715                   & 29017                  & 165,9    \\ \hline
8715                   & 29017                  & 165,9    \\ \hline
\end{tabular}
\caption{Зависимость времени полной верификации от сложности блока}
\label{tbl:models}
\end{table}

При сложности ... время достигло таймаута, и доказательство было остановлено до достижения однозначного ответа. По имеющимся данным, мы видим зависимость, похожую на кубическую. 

Для второй подзадачи выберем 
